\section{Glueing functions}
\label{section-tilde-module-sheaf}

\noindent
In this section we show that given an open covering
$$
\Spec(R) = \bigcup\nolimits_{i = 1}^n D(f_i)
$$
by standard opens, and given an element $h_i \in R_{f_i}$
for each $i$ such that $h_i = h_j$ as elements of $R_{f_i f_j}$
then there exists a unique $h \in R$ such that the image of
$h$ in $R_{f_i}$ is $h_i$. This result can be interpreted
in two ways:
\begin{enumerate}
\item The rule $D(f) \mapsto R_f$ is a sheaf of rings
on the standard opens, see Sheaves, Section \ref{sheaves-section-bases}.
\item If we think of elements of $R_f$ as the ``algebraic''
or ``regular'' functions on $D(f)$, then these glue
as would continuous, resp.\ differentiable functions
on a topological, resp.\ differentiable manifold.
\end{enumerate}

\begin{lemma}
\label{lemma-cover-module}
Let $R$ be a ring. Let $f_1, \ldots, f_n$ be elements of $R$
generating the unit ideal. Let $M$ be an $R$-module.
The sequence
$$
0 \to
M \xrightarrow{\alpha}
\bigoplus\nolimits_{i = 1}^n M_{f_i} \xrightarrow{\beta}
\bigoplus\nolimits_{i, j = 1}^n M_{f_i f_j}
$$
is exact, where $\alpha(m) = (m/1, \ldots, m/1)$
and $\beta(m_1/f_1^{e_1}, \ldots, m_n/f_n^{e_n})
= (m_i/f_i^{e_i} - m_j/f_j^{e_j})_{(i, j)}$.
\end{lemma}

\begin{proof}
It suffices to show that the localization of the sequence at
any maximal ideal $\mathfrak m$ is exact, see
Lemma \ref{lemma-characterize-zero-local}.
Since $f_1, \ldots, f_n$ generate the unit ideal,
there is an $i$ such that $f_i \not \in \mathfrak m$.
After renumbering we may assume $i = 1$.
Note that $(M_{f_i})_\mathfrak m = (M_\mathfrak m)_{f_i}$
and $(M_{f_if_j})_\mathfrak m = (M_\mathfrak m)_{f_if_j}$, see
Proposition \ref{proposition-localize-twice-module}.
In particular $(M_{f_1})_\mathfrak m = M_\mathfrak m$ and
$(M_{f_1 f_i})_\mathfrak m = (M_\mathfrak m)_{f_i}$, because
$f_1$ is a unit.
Note that the maps in the sequence are the canonical ones
coming from
Lemma \ref{lemma-universal-property-localization-module}
and the identity map on $M$.
Having said all of this, after replacing $R$ by $R_\mathfrak m$,
$M$ by $M_\mathfrak m$, and $f_i$ by their image in $R_\mathfrak m$,
and $f_1$ by $1 \in R_\mathfrak m$,
we reduce to the case where $f_1 = 1$.

\medskip\noindent
Assume $f_1 = 1$. Injectivity of $\alpha$ is now trivial. Let
$m = (m_i) \in \bigoplus_{i = 1}^n M_{f_i}$ be in the kernel of $\beta$.
Then $m_1 \in M_{f_1} = M$. Moreover, $\beta(m) = 0$
implies that $m_1$ and $m_i$ map to the same element of
$M_{f_1f_i} = M_{f_i}$. Thus $\alpha(m_1) = m$ and the
proof is complete.
\end{proof}

\begin{lemma}
\label{lemma-standard-covering}
Let $R$ be a ring, and let $f_1, f_2, \ldots f_n\in R$ generate
the unit ideal in $R$.
Then the following sequence is exact:
$$
0 \longrightarrow
R \longrightarrow
\bigoplus\nolimits_i R_{f_i} \longrightarrow
\bigoplus\nolimits_{i, j}R_{f_if_j}
$$
where the maps $\alpha : R \longrightarrow \bigoplus_i R_{f_i}$
and $\beta : \bigoplus_i R_{f_i} \longrightarrow \bigoplus_{i, j} R_{f_if_j}$
are defined as
$$
\alpha(x) = \left(\frac{x}{1}, \ldots, \frac{x}{1}\right)
\text{ and }
\beta\left(\frac{x_1}{f_1^{r_1}}, \ldots, \frac{x_n}{f_n^{r_n}}\right)
=
\left(\frac{x_i}{f_i^{r_i}}-\frac{x_j}{f_j^{r_j}}~\text{in}~R_{f_if_j}\right).
$$
\end{lemma}

\begin{proof}
Special case of Lemma \ref{lemma-cover-module}.
\end{proof}

\noindent
The following we have already seen above, but we state it explicitly here
for convenience.

\begin{lemma}
\label{lemma-disjoint-implies-product}
Let $R$ be a ring.
If $\Spec(R) = U \amalg V$ with both $U$ and $V$ open
then $R \cong R_1 \times R_2$ with $U \cong \Spec(R_1)$
and $V \cong \Spec(R_2)$ via the maps in Lemma \ref{lemma-spec-product}.
Moreover, both $R_1$ and $R_2$ are localizations as well as quotients
of the ring $R$.
\end{lemma}

\begin{proof}
By Lemma \ref{lemma-disjoint-decomposition} we have
$U = D(e)$ and $V = D(1-e)$ for some idempotent $e$.
By Lemma \ref{lemma-standard-covering} we see that
$R \cong R_e \times R_{1 - e}$ (since clearly $R_{e(1-e)} = 0$
so the glueing condition is trivial; of course it is
trivial to prove the product decomposition directly in this
case). The lemma follows.
\end{proof}

\begin{lemma}
\label{lemma-when-injective-covering}
Let $R$ be a ring.
Let $f_1, \ldots, f_n \in R$.
Let $M$ be an $R$-module.
Then $M \to \bigoplus M_{f_i}$ is injective if and only if
$$
M \longrightarrow \bigoplus\nolimits_{i = 1, \ldots, n} M, \quad
m \longmapsto (f_1m, \ldots, f_nm)
$$
is injective.
\end{lemma}

\begin{proof}
The map $M \to \bigoplus M_{f_i}$ is injective if and only if
for all $m \in M$ and $e_1, \ldots, e_n \geq 1$ such that
$f_i^{e_i}m = 0$, $i = 1, \ldots, n$ we have $m = 0$.
This clearly implies the displayed map is injective.
Conversely, suppose the displayed map is injective and
$m \in M$ and $e_1, \ldots, e_n \geq 1$ are such that
$f_i^{e_i}m = 0$, $i = 1, \ldots, n$. If $e_i = 1$ for all $i$,
then we immediately conclude that $m = 0$ from the injectivity of
the displayed map. Next, we prove this holds for any such data
by induction on $e = \sum e_i$. The base case is $e = n$, and we have
just dealt with this. If some $e_i > 1$, then set $m' = f_im$.
By induction we see that $m' = 0$. Hence we see that $f_i m = 0$,
i.e., we may take $e_i = 1$ which decreases $e$ and we win.
\end{proof}

\noindent
The following lemma is better stated and proved in the more general
context of flat descent. However, it makes sense to state it here
since it fits well with the above.

\begin{lemma}
\label{lemma-glue-modules}
Let $R$ be a ring. Let $f_1, \ldots, f_n \in R$. Suppose we are given
the following data:
\begin{enumerate}
\item For each $i$ an $R_{f_i}$-module $M_i$.
\item For each pair $i, j$ an $R_{f_if_j}$-module isomorphism
$\psi_{ij} : (M_i)_{f_j} \to (M_j)_{f_i}$.
\end{enumerate}
which satisfy the ``cocycle condition'' that all the diagrams
$$
\xymatrix{
(M_i)_{f_jf_k}
\ar[rd]_{\psi_{ij}}
\ar[rr]^{\psi_{ik}}
& &
(M_k)_{f_if_j} \\
&
(M_j)_{f_if_k} \ar[ru]_{\psi_{jk}}
}
$$
commute (for all triples $i, j, k$). Given this data define
$$
M = \Ker\left(
\bigoplus\nolimits_{1 \leq i \leq n} M_i
\longrightarrow
\bigoplus\nolimits_{1 \leq i, j \leq n} (M_i)_{f_j}
\right)
$$
where $(m_1, \ldots, m_n)$ maps to the element whose
$(i, j)$th entry is $m_i/1 - \psi_{ji}(m_j/1)$.
Then the natural map $M \to M_i$ induces an isomorphism
$M_{f_i} \to M_i$. Moreover $\psi_{ij}(m/1) = m/1$
for all $m \in M$ (with obvious notation).
\end{lemma}

\begin{proof}
To show that $M_{f_1} \to M_1$ is an isomorphism, it suffices
to show that its localization at every prime $\mathfrak p'$
of $R_{f_1}$ is an isomorphism, see
Lemma \ref{lemma-characterize-zero-local}.
Write $\mathfrak p' = \mathfrak p R_{f_1}$
for some prime $\mathfrak p \subset R$, $f_1 \not \in \mathfrak p$, see
Lemma \ref{lemma-standard-open}.
Since localization is exact
(Proposition \ref{proposition-localization-exact}),
we see that
\begin{align*}
(M_{f_1})_{\mathfrak p'} & =
M_\mathfrak p \\
& =
\Ker\left(
\bigoplus\nolimits_{1 \leq i \leq n} M_{i, \mathfrak p}
\longrightarrow
\bigoplus\nolimits_{1 \leq i, j \leq n} ((M_i)_{f_j})_\mathfrak p
\right) \\
& =
\Ker\left(
\bigoplus\nolimits_{1 \leq i \leq n} M_{i, \mathfrak p}
\longrightarrow
\bigoplus\nolimits_{1 \leq i, j \leq n} (M_{i, \mathfrak p})_{f_j}
\right)
\end{align*}
Here we also used Proposition \ref{proposition-localize-twice-module}.
Since $f_1$ is a unit in $R_\mathfrak p$, this reduces us to the case
where $f_1 = 1$ by replacing $R$ by $R_\mathfrak p$, $f_i$ by the
image of $f_i$ in $R_\mathfrak p$, $M$ by $M_\mathfrak p$, and
$f_1$ by $1$.

\medskip\noindent
Assume $f_1 = 1$. Then $\psi_{1j} : (M_1)_{f_j} \to M_j$
is an isomorphism for $j = 2, \ldots, n$. If we use these
isomorphisms to identify $M_j = (M_1)_{f_j}$, then we see
that $\psi_{ij} : (M_1)_{f_if_j} \to (M_1)_{f_if_j}$ is
the canonical identification. Thus the complex
$$
0 \to M_1 \to
\bigoplus\nolimits_{1 \leq i \leq n} (M_1)_{f_i}
\longrightarrow
\bigoplus\nolimits_{1 \leq i, j \leq n}
(M_1)_{f_if_j}
$$
is exact by Lemma \ref{lemma-cover-module}.
Thus the first map identifies $M_1$ with $M$ in this case
and everything is clear.
\end{proof}
